\chapter{BASIC EDITING}%This should be in all capital letters, "all caps!"

\section{Using this template}

To create your chapter 1, find the file named ``01Chapter.tex'' and open it in a latex editor. 
Then, replace all of this text. 
This file cannot be compiled directly; compiling this text and figures into a pdf is achieved by compiling the `OUThesis.tex' file.
Please try to compile a duplicate of this pdf before starting. 
Pay particular attention to making the ``List of Abbreviations'' in \autoref{AcronymSetion}

\section{Text Editing}

In Microsoft Word, typing occurs in paragraphs. 
For proof reading, locating the mistakes in their paragraph and correcting them in Microsoft Word is relatively easy. 
In LaTeX, however, correcting the mistakes during proof reading can at times a little difficult. 
To make it easier, it is suggested to write each sentence on its own line. 
This way, when there is an error to correct, one can simply look at the first word of the sentence to locate the text that should be changed.

\section{The OUThesis.tex File}

As stated earlier, to compile your dissertation, the `OUThesis.tex' file should be compiled. 
When starting, there are a number of things that should be adjusted in that file when work is begun.
There are comments in the OUThesis.tex file to explain this. 

The preamble is the portion of the file from top to `\textbackslash begin{document}'. The following should be adjusted in the preamble:
\begin{itemize}
	\item documentclass: If you are a PhD student writing your dissertation, the document class does not need to be changed.
If you are a masters student, the document class should be changed to read `\textbackslash documentclass[thesis, bookmarks]\{grizz\_thesis\}'.

	\item includeonly: The `\textbackslash includeonly{}' command can be uncommented to compile a subset of your document.

	\item usepackage: If you wish to use additional packages, they can be added to the preamble. 
The packages already in use can be found in the `grizz\_thesis.cls' file.

	\item custom macros: For a long thesis, it may be a good idea to use macros to make typing it in easier. 
For example, $\om$ and $\vst$ have macros defined.

	\item newacronym: The `\textbackslash newacronym' part of the preamble will be discussed in \autoref{AcronymSetion}.

	\item year: The year in your document can be changed in the preamble.
\end{itemize}

Between `\textbackslash begin\{document\}' and `\textbackslash end\{document\}' in the `OUThesis.tex' file, there are some things that need to be changed. 
The main thing is to enter the title, name, degree, and committee names. 
Besides this, the additional chapters and appendices should be uncommented as your thesis grows.



\section{Equations}

The beauty of latex is that if is excellent for formatting equations, and once you have the syntax memorized, equations are fast to type. 
For example, the equation $\alpha(x) = \beta\delta y\sqrt{\epsilon}$ is typed inline between two \$ signs. 
If instead an equation should be numbered, it might be best to code it as:
\begin{equation} \alpha(x) = \beta\delta y\sqrt{\epsilon} \label{first}\end{equation}
To have an equation without numbering, simply use the following:
\[\alpha(x) = \beta\delta y\sqrt{\epsilon}\]
%This is how equations are entered.

How, one might ask, should an equation be referenced in the text? 
The easy way is to use the included package, ``hyperref''. 
This allows you to refer to \autoref{first} and \autoref{second} and \autoref{third}.
As will be shown in the next section, ``\textbackslash autoref\{\}'' also works with figures and tables and the like. 
Please note that despite their black color, hyperref generates hyperlinks in the pdf output. 
Thus, if you click on this: \autoref{first} your pdf viewer might take you back a page. 

If you want a piecewise equation with a single number, here is an example.
\begin{equation}
 r(t) = \begin{cases} A \cos(\pi \rho t^2 + \om t + \psi_r) &-\tfrac{T}{2}\leq t\leq\tfrac{T}{2} \\
 0 &\text{otherwise}\end{cases} \label{sr1}
\end{equation}


If instead the pieces should be referenced individually, try:

\begin{subequations}
\label{mainVOltage}
\begin{numcases}{v(t) = }
-\tfrac{1}{2} I\Delta R \cos( \psi) & when injection locked ~~~~~~~~\label{second}\\
-\tfrac{1}{2} I\Delta R \cos(\pi \rho t^2 + \om t + \psi) & otherwise \label{third}
\end{numcases}
\end{subequations}

\section{Quotations}

When you have a long quotation, it may require a quotation block. 
Please refer to the OU thesis formatting guidelines to determine when a block is required. 

\begin{quote}
It was actually Bart Cameron's error and you'll have to understand about Bart Cameron. 
He's the sheriff at Twin Gulch, Idaho, and I'm his deputy. 
Bart Cameron is an impatient man and he gets most impatient when he has to work up his income tax. 
You see, besides being sheriff, he also owns and runs the general store, he's got some shares in a sheep ranch, he does a bit of assay work, he's got a kind of pension for being a disabled veteran (bad knee) and a few other things like that. 
Naturally, it makes his tax figures complicated.
\end{quote}
This quote was taken from an Issav Assimov story, \textit{The Dead Past} \cite{asimov1956dead}. 


\section{Acronyms\label{AcronymSetion}}

In the front matter, there is a page called ``List of Abbreviations''. 
How was this generated? 

First, it is important to discuss the `glossaries' package.
The glossaries package is included in the grizz\_thesis.cls file. 
This allows the use of the \textbackslash gls\{\} and \textbackslash glspl\{\} commands. 
In the preamble of the `OUThesis.tex' file, there is an line \textbackslash newacronym\{fpga\}\{FPGA\}\{Field-Programmable Gate Array\}. 
This can be referenced in singular with the code \textbackslash gls\{fpga \}, and with this command it will display as \gls{fpga}. 
As can be seen, the first time \textbackslash gls\{fpga \} is used it displays the full name and the acronym in parenthesis. 
The second time this command is used, it will just say \gls{fpga}.
If is should be plural, it must use the code \textbackslash glspl\{\}. 
Then, it will be displayed as \glspl{fpga}.

Once the glossaries package is being used correctly, the List of Abbreviations can be generated. 
This is a little complicated - it requires that perl be installed on your computer. 
If you are using Windows, add the directory of perl to the path. 
After this, run the command `makeglossaries OUThesis' inside the directory of your latex files. 
This will generate the files required to populate the list of abbreviations.
If you are using windows, there is a batch file `makeGloss.bat' that can be used with a mouse in Windows.
The command `makeglossaries.exe' should also be in the path. 
As I used MiKTeX, it was found in the MiKTeX directory.
There may be some trouble shooting associated with this. 

Once all these steps are done, the List of Abbreviations should appear correctly in the generated pdf file.

%\begin{figure}[tbp]
%\centering
%\includegraphics{figures/NayaritMexico}
%\caption[First Chapter Figure]{First Chapter Figure. PDF vector image is six inches wide, and copied from Wikipedia. Figures are discussed in Chapter 2.}
%\label{firstChapterFigureLabel}
%\end{figure}